\documentclass[10pt,a4paper]{amsart}
\usepackage[margin=19mm]{geometry}
\usepackage{amsmath,amssymb,mathtools}
\usepackage[scr=rsfs,cal=euler]{mathalfa}
\usepackage{xfrac}
\usepackage[unicode,hidelinks,bookmarks=false]{hyperref}
\newcommand{\ie}{i.e.\ }
\newcommand{\cf}{cf.\ }
\newcommand{\resp}{respectively\ }
\newcommand{\Subsection}{\subsection}
\providecommand{\nobreakdash}{\nobreak\hbox{-}}
\setlength{\emergencystretch}{2em}
\multlinegap0pt
\usepackage{bm}
\usepackage{bbm}
\usepackage{dsfont}
\usepackage{xspace}
\usepackage{cancel}
\usepackage{mathtools}
\usepackage{enumitem}
\usepackage{amsthm}
\makeatletter
\@ifundefined{theorem}{\newtheorem{theorem}{Theorem}}{}
\@ifundefined{definition}{\newtheorem{definition}[theorem]{Definition}}{}
\@ifundefined{lemma}{\newtheorem{lemma}[theorem]{Lemma}}{}
\@ifundefined{proposition}{\newtheorem{proposition}[theorem]{Proposition}}{}
\makeatother
\DeclareMathOperator{\hol}{Hol}
\newcommand{\Ra}{\Rightarrow}
\newcommand{\bep}[1]{\bm{\varepsilon_{#1}}}
\newcommand{\mC}{\mathbb{C}}
\newcommand{\mM}{\mathbb{M}}
\newcommand{\ra}{\rightarrow}
\title[Classes of Holomorphic Multicomplex-Valued Functions Generat]{Classes of Holomorphic Multicomplex-Valued Functions Generated by Elliptic-Admissible Involutions}
\author{Nicolas Doyon, Pierre-Olivier Paris\'e, William Verreault}
\date{Source: arXiv:2211.13875v2 (CC BY 4.0)}
\begin{document}
\maketitle

\noindent\emph{Source and licence.} Classes of Holomorphic Multicomplex-Valued Functions Generated by Elliptic-Admissible Involutions. Version arXiv:2211.13875v2 (CC BY 4.0), distributed under the Creative Commons Attribution 4.0 International licence, as recorded in the arXiv licence metadata for this exact version and shown on the abstract page https://arxiv.org/abs/2211.13875v2. Full rights notice: licenses/arxiv-2211.13875-CC-BY-4.0.txt.\par\smallskip

\noindent\textit{Editorial scope.} Counted unit: the Proposition whose source label is \texttt{P:DerivativeEquivalence} in the article (source0.tex lines 430--460, source label \texttt{P:DerivativeEquivalence}), delivered together with the article's own complete original proof of it (source0.tex lines 462--616, one \texttt{proof} environment of 155 lines). No mathematics is written, repaired or re-derived: statement and proof are the article's own text line for line. Required context retained uncounted: L:Invertibility (lines 389--397); D:MCHolo (lines 411--417). Every definition, display and label of those lines is retained unchanged. Omitted from this package and not redistributed: 1--388, 398--410, 418--429, 461--461, 617--1662. Nothing inside a retained excerpt is omitted or altered: every retained body is delivered exactly as the article's lines are written, so no omission receipt of any kind accompanies this package. Citations: the retained excerpts carry no citation command at all, so no bibliography entry of the article is transcribed and no reference is redistributed. Reference adaptation: none was needed and none was made; every reference inside the delivered unit resolves inside it, so no unresolved reference or citation is produced. Printed slips kept literal: no printed word, symbol, hypothesis or number of the counted statement or of its proof was altered. Layout and class: the article's own class is \texttt{amsart} with packages including lmodern, amssymb, amsmath, amsthm, enumerate, enumitem, geometry, xcolor, multicol, hyperref, comment, cleveref, bm, bbm; the wrapper is the lane's pinned \texttt{amsart} editorial wrapper at 10pt on A4, which restates the theorem-like environments the retained text needs and adds the article's own macro definitions that the retained text needs (\texttt{\textbackslash Ra}, \texttt{\textbackslash bep}, \texttt{\textbackslash hol}, \texttt{\textbackslash mC}, \texttt{\textbackslash mM}, \texttt{\textbackslash ra}), with extra wrapper packages (bm, bbm, dsfont, xspace, cancel, mathtools, enumitem). The delivered numbering is the unit's own. Assets: no retained span contains a figure environment, a graphics-inclusion command, a PDF-page inclusion, a table or a TikZ picture; no image file is copied into this package, the package declares no asset dependency and no image file is redistributed with it. Licence: version arXiv:2211.13875v2 of this article is distributed under the Creative Commons Attribution 4.0 International licence, as recorded in the arXiv licence metadata for this exact version and shown on the abstract page https://arxiv.org/abs/2211.13875v2. A full rights notice is delivered at licenses/arxiv-2211.13875-CC-BY-4.0.txt. Characters: every retained excerpt is pure ASCII, so the delivered engine, XeLaTeX with its default Latin Modern text font, reports no missing character.
    \begin{lemma}\label{L:Invertibility}
    A multicomplex number $\eta \in \mM \mC(n)$ is invertible if and only if
    $\eta_{\bep{}} \neq 0$ for every $\bep{} \in \mathcal{E}_n$, in which case
        \[
        \eta^{-1} = \sum_{\bep{} \in \mathcal{E}_n} \eta_{\bep{}}^{-1}\, \bep{} .
        \]
    Moreover, the group $\mM \mC(n)^\times$ of invertible elements is open and
    dense in $\mM \mC(n)$.
    \end{lemma}

    \begin{definition}\label{D:MCHolo}
    A function $F : U \ra \mM \mC(n)$ of class $C^1$ is \emph{multicomplex holomorphic} on $U$ if, for every $\eta \in U$, the differential $D F (\eta)$ is $\mM \mC(n)$-linear, that is, if there exists $\lambda (\eta ) \in \mM \mC(n)$ such that
        \begin{align*}
        DF (\eta ) [h] = \lambda (\eta )\, h , \qquad h \in \mM \mC(n).
        \end{align*}
    We denote by $\hol(U)$ the set of multicomplex holomorphic functions on $U$.
    \end{definition}

    \begin{proposition}\label{P:DerivativeEquivalence}
    Let $U \subseteq \mM \mC(n)$ be open and let $F : U \ra \mM \mC(n)$ be a
    function. The following statements are equivalent.
        \begin{enumerate}[label=(\roman*)]
        \item $F$ is multicomplex holomorphic on $U$ in the sense of
        Definition~\ref{D:MCHolo}.
        \item For every $\eta_0 \in U$, the limit
            \begin{equation}\label{Eq:MCDerivative}
            F' (\eta_0 ) := \lim_{\substack{h \ra 0 \\ h \in \mM \mC(n)^\times}}
            \frac{ F (\eta_0 + h ) - F (\eta_0 )}{h} 
            \end{equation}
        exists in $\mM \mC(n)$, where $h$ ranges over the invertible elements
        such that $\eta_0 + h \in U$.
        \item Every point of $U$ admits an idempotent polydisc neighborhood
        $P ( a , r ) \subseteq U$ on which
            \begin{equation}\label{Eq:LocalSplitting}
            F (\eta ) = \sum_{\bep{} \in \mathcal{E}_n } \varphi_{\bep{}} ( \eta_{\bep{}} )\, \bep{} ,
            \qquad \eta \in P ( a , r ),
            \end{equation}
        where each $\varphi_{\bep{}}$ is a holomorphic function of one complex
        variable on the disc
        $D_{\bep{}} := \{ z \in \mM \mC(1) : | z - a_{\bep{}} | < r \}$.
        \end{enumerate}
    In this case, $F$ is of class $C^\infty$ on $U$ and, for every
    $\eta \in U$,
        \begin{equation}\label{Eq:DerivativeIdentification}
        F' (\eta ) = \lambda (\eta ) = DF (\eta ) [ 1 ]
        = \sum_{\bep{} \in \mathcal{E}_n} \varphi_{\bep{}} ' ( \eta_{\bep{}} )\, \bep{}
        \end{equation}
    in the local representation \eqref{Eq:LocalSplitting}.
    \end{proposition}

    \begin{proof}
    We prove that (iii) implies (i) and (ii), that (i) implies (iii), and
    finally that (ii) implies (iii). The identities in
    \eqref{Eq:DerivativeIdentification} are established along the way.

    (iii) $\Ra$ (i). Suppose that \eqref{Eq:LocalSplitting} holds on
    $P = P ( a , r ) \subseteq U$. 
    For each $\bep{} \in \mathcal{E}_n$, define $\pi_{\bep{}} : \mM \mC (n) \ra \mM \mC (1)$ by $\pi_{\bep{}} (\eta) := \eta_{\bep{}}$ for $\eta = \sum_{\bep{} \in \mathcal{E}_n} \eta_{\bep{}} \bep{}$ and $\iota_{\bep{}} : \mM \mC (1) \ra \mM \mC (n)$ by $\iota_{\bep{}} (z) := z \bep{}$, so that $F = \sum_{\bep{} \in \mathcal{E}_n} G_{\bep{}}$ on $P$, where $G_{\bep{}} := \iota_{\bep{}} \circ \varphi_{\bep{}} \circ \pi_{\bep{}}$. 
    Each map $\pi_{\bep{}}$ and $\iota_{\bep{}}$ is real-linear, and each $\varphi_{\bep{}}$ is holomorphic, hence of class $C^\infty$. Consequently $F$ is of class $C^\infty$ on $P$ and therefore on $U$, since $U$ is covered by such polydiscs. 
    Moreover, by the chain rule and the complex-linear nature of the differential of a holomorphic complex-valued function, 
    $$
        DG_{\bep{}} (\eta) [h] = \iota_{\bep{}} [ D\varphi_{\bep{}} (\eta_{\bep{}}) [ \pi_{\bep{}} [h]]] = \varphi_{\bep{}}' (\eta_{\bep{}}) h_{\bep{}} \bep{} ,
    $$
    and therefore
        \[
        DF (\eta ) [ h ]
        = \sum_{\bep{} \in \mathcal{E}_n}\varphi_{\bep{}} ' ( \eta_{\bep{}} )\, h_{\bep{}}\, \bep{}
        = \Big( \sum_{\bep{} \in \mathcal{E}_n} \varphi_{\bep{}} ' ( \eta_{\bep{}} )\, \bep{} \Big)\, h ,
        \qquad h \in \mM \mC(n) ,
        \]
    the second equality is because multiplication is componentwise. Thus
    $DF (\eta )$ is $\mM \mC(n)$-linear with
    $\lambda (\eta ) = \sum_{\bep{} \in \mathcal{E}_n} \varphi_{\bep{}} ' ( \eta_{\bep{}} ) \bep{}$, so
    $F \in \hol (U)$, and evaluating at $h = 1$ gives
    $\lambda (\eta ) = DF (\eta ) [ 1 ]$.

    (iii) $\Ra$ (ii). Fix $\eta_0 \in U$. Choose the polydisc in (iii) centered at $a = \eta_0$. Let $h \in \mM \mC(n)^\times$
    with $\eta_0 + h \in P ( a , r )$. By Lemma~\ref{L:Invertibility},
    $h_{\bep{}} \neq 0$ for every $\bep{}$ and inversion is componentwise, so that
        \[
        \frac{F (\eta_0 + h ) - F (\eta_0 )}{h} 
        = \sum_{\bep{} \in \mathcal{E}_n}
        \frac{ \varphi_{\bep{}} ( a_{\bep{}} + h_{\bep{}} ) - \varphi_{\bep{}} ( a_{\bep{}} ) }{ h_{\bep{}} }\,
        \bep{} .
        \]
    As $h \ra 0$, each component $h_{\bep{}}$ tends to $0$ while remaining nonzero,
    and since $\varphi_{\bep{}}$ is complex-differentiable at $a_{\bep{}}$, the ${\bep{}}$-th
    component of the right-hand side converges to $\varphi_{\bep{}} ' ( a_{\bep{}} )$.
    componentwise convergence being equivalent to convergence in
    $\mM \mC(n)$, the limit \eqref{Eq:MCDerivative} exists and equals
    $\sum_{\bep{} \in \mathcal{E}_n} \varphi_{\bep{}} ' ( a_{\bep{}} ) \bep{} = F'(\eta_0)$. This proves
    (ii).

    (i) $\Ra$ (iii). Assume that $F$ is of class $C^1$ on $U$ and that $DF (\eta )$ is multicomplex linear for any $\eta \in U$. Let $a \in U$ and choose $r > 0$ with $P := P ( a , r ) \subseteq U$. Writing $F = \sum_{\bep{} \in \mathcal{E}_n} F_{\bep{}} \bep{}$ and $\lambda = \sum_{\bep{} \in \mathcal{E}_n} \lambda_{\bep{}} \bep{}$, we see that $F_{\bep{}} = \pi_{\bep{}} \circ F$ and $\lambda_{\bep{}} = \pi_{\bep{}} \circ \lambda$. Therefore, from the chain rule, we obtain
        \begin{equation}\label{Eq:ComponentwiseDifferential}
        D F_{\bep{}} (\eta ) [ h ] = \pi_{\bep{}} [DF (\eta ) [h] ] = \pi_{\bep{}} [\lambda (\eta ) h] = \lambda_{\bep{}} (\eta )\, h_{\bep{}} ,
        \end{equation}
    for $\eta \in P$, $h \in \mM \mC(n)$, and $ \bep{} \in \mathcal{E}_n$. 

    Fix $\bep{}$ and let $\bep{}' \neq \bep{}$. If $\eta , \zeta \in P$ differ only in
    their $\bep{}'$ components, the segment joining them lies in $P$ and is of
    the form $t \mapsto \eta + t\, w\, \bep{}'$, $t \in [ 0 , 1 ]$, with $w = \zeta_{\bep{}'} - \eta_{\bep{}'} \in \mM \mC (1)$.
    By \eqref{Eq:ComponentwiseDifferential}, the derivative of the $C^1$ map $t \mapsto F_{\bep{}} ( \eta + t\, w\, \bep{}' )$
    equals $D F_{\bep{}} ( \cdot ) [ w \bep{}' ] = \lambda_{\bep{}} ( \cdot )\,
    \pi_{\bep{}} ( w \bep{}' ) = 0$, since $\pi_{\bep{}} ( w \bep{}' ) = 0$ for $\bep{} \neq \bep{}'$. Hence
    $F_{\bep{}} ( \zeta ) = F_{\bep{}} ( \eta )$ by the mean value theorem applied to the real components of $F_{\bep{}}$. Since any two points of $P$ with the
    same $\bep{}$ component differ by finitely many such one-component moves
    inside $P$, the component $F_{\bep{}}$ depends only on $\eta_{\bep{}}$ on $P$. Therefore there is
    a function $\varphi_{\bep{}} : D_{\bep{}} \ra \mM \mC(1)$ with
    $F_{\bep{}} (\eta ) = \varphi_{\bep{}} ( \eta_{\bep{}} )$ for $\eta \in P$.

    It remains to check that each $\varphi_{\bep{}}$ is holomorphic on $D_{\bep{}}$. Let
    $z \in D_{\bep{}}$ and set $\eta := a + ( z - a_{\bep{}} )\, \bep{} \in P$. The map
    $z ' \mapsto \varphi_{\bep{}} ( z ' ) = F_{\bep{}} \big( \eta + ( z ' - z ) \bep{} \big)$
    is the composition of a real-affine map with $F_{\bep{}}$, hence is differentiable at $z$, and by \eqref{Eq:ComponentwiseDifferential} its
    differential is $w \mapsto D F_{\bep{}} (\eta ) [ w \bep{} ] =
    \lambda_{\bep{}} (\eta )\, w$, which is $\mM \mC(1)$-linear. Thus $\varphi_{\bep{}}$ is
    complex-differentiable at every point of $D_{\bep{}}$, hence holomorphic on
    $D_{\bep{}}$. This proves (iii).

    (ii) $\Ra$ (iii). This is the only implication requiring some care, since
    the increments in \eqref{Eq:MCDerivative} must avoid the null cone and no
    regularity of $F$ is assumed. Let $a \in U$ and choose $r > 0$ so that
    $P := P ( a , r ) \subseteq U$.

    We first record the estimate provided by (ii). Let $q \in P$. Applying
    \eqref{Eq:MCDerivative} at $q$, there exists
    $\delta_q > 0$, which we may take small enough that
    $\| h \| < \delta_q$ implies $q + h \in U$, such that for $h \in \mM \mC(n)^\times$, $0 < \| h \| < \delta_q$,
        \begin{equation}\label{Eq:BasicEstimate}
        \big\| \big( F ( q + h ) - F ( q ) \big)\, h^{-1} - F ' ( q ) \big\|
        \leq 1 .
        \end{equation}
    For such $h$, write $Q := ( F ( q + h ) - F ( q ) )\, h^{-1}$, so that
    $F ( q + h ) - F ( q ) = Q\, h$ exactly. Taking the $\bep{}$ component and
    using $| Q_{\bep{}} | \leq \| F ' ( q ) \| + 1$, which follows from
    \eqref{Eq:BasicEstimate}, we deduce
        \begin{equation}\label{Eq:ComponentEstimate}
        \big| F_{\bep{}} ( q + h ) - F_{\bep{}} ( q ) \big| \leq C_q\, | h_{\bep{}} | ,
        \qquad C_q := \| F ' ( q ) \| + 1 , \quad \bep{} \in \mathcal{E}_n.
        \end{equation}
    We now show our main goal in two steps: (1) we first show that, on $P$, each component $F_{\bep{}}$ depends only on $\eta_{\bep{}}$; and (2) we then show that each $\varphi_{\bep{}}$ is holomorphic on $D_{\bep{}}$, which will complete the proof of the equivalence.

    (1) Fix $\bep{} \neq \bep{}'$ and let $\eta , \zeta \in P$ differ only in their
    $\bep{}'$ component, say $\zeta = \eta + w\, \bep{}'$ with
    $w \in \mM \mC(1) \setminus \{ 0 \}$. For $t \in [ 0 , 1 ]$, set
    $q_t := \eta + t\, w\, \bep{}' \in P$, and abbreviate
    $\delta_t := \delta_{q_t}$ and $C_t := C_{q_t}$. We claim that
    $t \mapsto F_{\bep{}} ( q_t )$ is locally constant on $[ 0 , 1 ]$, meaning that there is a neighborhood around any point in $[0, 1]$ on which $t \mapsto F_{\bep{}} (q_t)$ is constant.

    Fix $t \in [ 0 , 1 ]$ and let $t ' \in [ 0 , 1 ]$ with
    $0 < | t ' - t |\, | w | < \delta_t / 2$. Let $s$ be a real number with
        \[
        0 < s < \min \big( \delta_t / 4 ,\ \delta_{t '} ,\
        | t ' - t |\, | w | \big) ,
        \]
    and consider the increment and the auxiliary point
        \[
        h := s + ( t ' - t )\, w\, \bep{}' , \qquad y := q_t + h .
        \]
    The components of $h$ are $h_{\bep{}''} = s$ for $\bep{}'' \neq \bep{}'$ and
    $h_{\bep{}'} = s + ( t ' - t )\, w$, and are nonzero by the choice of $s$. Hence $h \in \mM \mC(n)^\times$ by
    Lemma~\ref{L:Invertibility}, and
    $\| h \| \leq s + | t ' - t |\, | w | < \delta_t$. On the other hand,
    since $q_{t '} = q_t + ( t ' - t )\, w\, \bep{}'$, we also have
        \[
        y = q_{t '} + s , \qquad s \in \mM \mC(n)^\times , \quad
        \| s \| = s < \delta_{t '} ,
        \]
    the real number $s = s \cdot 1$ being invertible because all of its
    components equal $s \neq 0$. Applying \eqref{Eq:ComponentEstimate} at the
    point $q_t$ with the increment $h$, whose $\bep{}$ component is $s$ because
    $\bep{} \neq \bep{}'$, and then at the point $q_{t '}$ with the increment $s$, we
    obtain
        \[
        | F_{\bep{}} ( y ) - F_{\bep{}} ( q_t ) | \leq C_t\, s
        \qquad \text{and} \qquad
        | F_{\bep{}} ( y ) - F_{\bep{}} ( q_{t '} ) | \leq C_{t '}\, s .
        \]
    Hence $| F_{\bep{}} ( q_{t '} ) - F_{\bep{}} ( q_t ) | \leq ( C_t + C_{t '} )\, s$ for
    every sufficiently small $s > 0$, and therefore
    $F_{\bep{}} ( q_{t '} ) = F_{\bep{}} ( q_t )$. This holds in the neighborhood $I_t := (t - \rho_t , t + \rho_t) \cap [0, 1]$ where $\rho_t := \delta_t / (2 |w| )$. Hence, $t\mapsto F_{\bep{}}(q_{t'})$ is constant on $I_t$ and this proves what we wanted. Now, since $[0, 1]$ is a connected set and $t \mapsto F_{\bep{}} ( q_t )$ is locally constant, it must be constant on $[0, 1]$. Therefore, we obtain $F_{\bep{}} ( \zeta ) = F_{\bep{}} ( \eta )$. As in the proof of (i) $\Ra$ (iii), it follows that $F_{\bep{}}$ depends only on the component $\eta_{\bep{}}$ and hence $F_{\bep{}} ( \eta ) = \varphi_{\bep{}} ( \eta_{\bep{}} )$ in $P$ for some function $\varphi_{\bep{}} : D_{\bep{}} \ra \mM \mC(1)$.

    (2) Let $z \in D_{\bep{}}$
    and set $\eta := a + ( z - a_{\bep{}} )\, \bep{} \in P$. Let $( z_p )_{p \geq 1}$
    be any sequence in $D_{\bep{}} \setminus \{ z \}$ with $z_p \ra z$, and set
    $h_p := z_p - z \in \mM \mC(1) \setminus \{ 0 \} \subseteq \mM \mC (n)$. All the idempotent
    components of $h_p$ equal $z_p - z \neq 0$, so
    $h_p \in \mM \mC(n)^\times$ by Lemma~\ref{L:Invertibility}, while
    $\| h_p \| = | z_p - z | \ra 0$ and $\eta + h_p \in P$ for $p$ large
    enough. By (ii) applied at $\eta$, the quotients
    $( F ( \eta + h_p ) - F ( \eta ) )\, h_p^{-1}$ converge to $F ' (\eta )$. Applying $\pi_{\bep{}}$ and using Step (1), we obtain
        \[
        \frac{ \varphi_{\bep{}} ( z_p ) - \varphi_{\bep{}} ( z ) }{ z_p - z }
        \; \longrightarrow \; \pi_{\bep{}} \big( F ' ( \eta ) \big).
        \]
    Since the sequence $( z_p )$ was arbitrary, $\varphi_{\bep{}}$ is
    complex-differentiable at $z$, with
    $\varphi_{\bep{}} ' ( z ) = \pi_{\bep{}} \big( F ' ( \eta ) \big)$. Thus $\varphi_{\bep{}}$ is
    complex-differentiable at every point of the open disk $D_{\bep{}}$, hence
    holomorphic on $D_{\bep{}}$. This proves (iii).
    
    The identities in \eqref{Eq:DerivativeIdentification} were
    obtained along the way.
    \end{proof}

\end{document}
